\documentclass[a4paper,12pt]{article}

\usepackage[T1]{fontenc}
\usepackage[italian,english]{babel}
\usepackage[utf8]{inputenc}
\usepackage{graphicx}
\usepackage{float}
\usepackage{subcaption}
\usepackage{geometry}
\usepackage{fancyhdr}
\usepackage{xcolor}
\usepackage{mdframed}
\usepackage{enumitem}
\usepackage{hyperref}
\usepackage{booktabs}
\usepackage{array}
\usepackage{multicol}

% ── Colori brand ─────────────────────────────────────────────────────────────
\definecolor{nutrigreen}{RGB}{46,125,50}
\definecolor{nutrilightgreen}{RGB}{232,245,233}
\definecolor{nutridark}{RGB}{28,28,28}
\definecolor{nutrigray}{RGB}{117,117,117}
\definecolor{nutriaccent}{RGB}{56,142,60}

% ── Margini pagina ────────────────────────────────────────────────────────────
\setlength{\headheight}{30pt}
\geometry{left=2.5cm, right=2.5cm, top=2.8cm, bottom=2.5cm}

% ── Intestazione ─────────────────────────────────────────────────────────────
\pagestyle{fancy}
\fancyhf{}
\fancyhead[L]{\small\color{nutrigray}Sezione 4 --- Guida all'utilizzo dell'applicazione}
\fancyhead[R]{\small\color{nutrigray}\thepage}
\renewcommand{\headrulewidth}{0.4pt}
\renewcommand{\headrule}{\hbox to\headwidth{\color{nutrigreen}\leaders\hrule height \headrulewidth\hfill}}

% ── Contatore passi ───────────────────────────────────────────────────────────
\newcounter{passo}[subsection]
\newcommand{\passaggio}[1]{%
    \stepcounter{passo}%
    \subsubsection*{\color{nutriaccent}Passaggio \thesubsection.\thepasso\ ---\ #1}%
}

% ── Box informativo verde ─────────────────────────────────────────────────────
\newmdenv[
  backgroundcolor=nutrilightgreen,
  linecolor=nutrigreen,
  linewidth=1.5pt,
  roundcorner=5pt,
  innertopmargin=8pt,
  innerbottommargin=8pt,
  innerleftmargin=12pt,
  innerrightmargin=12pt,
  skipabove=8pt,
  skipbelow=8pt
]{infobox}

% ── Etichetta grassetto verde ─────────────────────────────────────────────────
\newcommand{\labelgreen}[1]{\medskip\noindent{\color{nutriaccent}\textbf{#1}}\par\smallskip}

% ── Suggerimento corsivo ──────────────────────────────────────────────────────
\newcommand{\suggerimento}[1]{%
  \smallskip\noindent\textit{\color{nutrigray}#1}\par\smallskip%
}

% ── Variabile per cartella immagini (utile per traduzioni) ───────────────────
\newcommand{\imgdir}{immagini_ita}

% ── Immagine centrata con didascalia ─────────────────────────────────────────
\newcommand{\schermata}[3][0.38\textwidth]{%
  \begin{figure}[H]
    \centering
    \includegraphics[width=#1]{\imgdir/#2}
    \caption{\small\color{nutrigray}#3}
  \end{figure}%
}

% ── Due schermate affiancate ──────────────────────────────────────────────────
\newcommand{\dueschermate}[4]{%
  \begin{figure}[H]
    \centering
    \begin{subfigure}[b]{0.42\textwidth}
      \centering
      \includegraphics[width=\textwidth]{\imgdir/#1}
      \caption{\small\color{nutrigray}#2}
    \end{subfigure}
    \hfill
    \begin{subfigure}[b]{0.42\textwidth}
      \centering
      \includegraphics[width=\textwidth]{\imgdir/#3}
      \caption{\small\color{nutrigray}#4}
    \end{subfigure}
  \end{figure}%
}

% ── Tre schermate affiancate ──────────────────────────────────────────────────
\newcommand{\treschermate}[6]{%
  \begin{figure}[H]
    \centering
    \begin{subfigure}[b]{0.30\textwidth}
      \centering
      \includegraphics[width=\textwidth]{\imgdir/#1}
      \caption{\small\color{nutrigray}#2}
    \end{subfigure}
    \hfill
    \begin{subfigure}[b]{0.30\textwidth}
      \centering
      \includegraphics[width=\textwidth]{\imgdir/#3}
      \caption{\small\color{nutrigray}#4}
    \end{subfigure}
    \hfill
    \begin{subfigure}[b]{0.30\textwidth}
      \centering
      \includegraphics[width=\textwidth]{\imgdir/#5}
      \caption{\small\color{nutrigray}#6}
    \end{subfigure}
  \end{figure}%
}

% ─────────────────────────────────────────────────────────────────────────────
\begin{document}

% ════════════════════════════════════════════════════════════════════════════
%  INTESTAZIONE SEZIONE
% ════════════════════════════════════════════════════════════════════════════
\begin{center}
  {\Huge\bfseries\color{nutrigreen} Sezione 4}\\[6pt]
  {\Large\bfseries\color{nutridark} Guida all'utilizzo dell'applicazione}\\[4pt]
  {\normalsize\itshape\color{nutrigray} Manuale utente --- istruzioni passo per passo}
\end{center}

\medskip
\noindent\rule{\textwidth}{1.5pt}{\color{nutrigreen}}
\medskip

\noindent Il presente capitolo illustra le funzionalita' principali di \textbf{NutriApp}
attraverso istruzioni dettagliate e schermate tratte direttamente dall'applicazione.
Ogni sezione e' suddivisa in passaggi progressivi per facilitare la comprensione anche
agli utenti meno esperti.

\medskip
\begin{infobox}
  \textbf{\color{nutrigreen}Indice delle sezioni}
  \begin{enumerate}[leftmargin=1.4em, itemsep=2pt]
    \item Accesso e registrazione
    \item Barra di navigazione
    \item Inserimento degli alimenti
    \item Impostazione degli obiettivi
    \item Visualizzazione del progresso
    \item Impostazioni dell'account
  \end{enumerate}
\end{infobox}

\newpage

% ════════════════════════════════════════════════════════════════════════════
%  1. ACCESSO E REGISTRAZIONE
% ════════════════════════════════════════════════════════════════════════════
\section{Accesso e registrazione}

Alla prima apertura dell'applicazione viene mostrata la schermata di benvenuto,
da cui e' possibile accedere a un account esistente oppure crearne uno nuovo.
NutriApp supporta tre modalita' di accesso: credenziali email e password,
autenticazione tramite account Google e autenticazione tramite Apple ID.
E' inoltre possibile esplorare l'app in modalita' ospite senza creare un account.

% ─── 1.1 ─────────────────────────────────────────────────────────────────────
\subsection{Schermata di benvenuto e login}

\passaggio{Aprire l'applicazione}

\schermata[0.40\textwidth]{login.jpeg}{Schermata di benvenuto \textbf{NutriApp}:
  in cima il logo verde e il sottotitolo \emph{Accedi per gestire i tuoi pasti,
  i macro e i micro nutrienti in un'unica app.} Il pannello bianco centrale
  mostra i campi \textbf{Email} e \textbf{Password} (con icona occhio per
  mostrare/nascondere i caratteri), il pulsante verde \textbf{ACCEDI} e il
  link \emph{Continua senza accedere}. Sotto il separatore \emph{oppure} i
  pulsanti \textbf{Accedi con Google} e \textbf{Accedi con Apple}, e in fondo
  il link \emph{Non hai un account? Registrati}.}

\labelgreen{Descrizione}
Alla prima apertura di NutriApp compare la schermata di accesso con il nome
dell'app in verde scuro e un sottotitolo che ne descrive brevemente la funzione.
Il pannello di accesso centrale presenta due campi:

\begin{itemize}[leftmargin=1.4em, itemsep=2pt]
  \item \textbf{Email} --- indirizzo di posta elettronica associato all'account,
        indicato dall'icona busta a sinistra del campo
  \item \textbf{Password} --- password dell'account; toccare l'icona occhio
        barrato a destra del campo per alternare la visualizzazione dei caratteri
        in chiaro o mascherati
\end{itemize}

Dopo aver compilato entrambi i campi, premere il pulsante verde \textbf{Accedi}
per completare l'autenticazione ed entrare nell'applicazione.

\labelgreen{Continua senza accedere}
Toccando il link \emph{Continua senza accedere} e' possibile esplorare
l'applicazione in modalita' ospite, senza creare ne' inserire credenziali.
In questa modalita' i dati non vengono sincronizzati sul cloud e non sono
accessibili da altri dispositivi. Si consiglia di creare un account per
non perdere i dati inseriti.

\suggerimento{Se si dimentica la password, e' possibile recuperarla tramite
la procedura di recupero disponibile sul sito ufficiale di NutriApp: il
sistema invia un link di reset all'indirizzo email associato all'account.}

% ─── 1.2 ─────────────────────────────────────────────────────────────────────
\subsection{Accedere con Google}

\passaggio{Selezionare l'accesso rapido tramite Google}

\dueschermate{login.jpeg}{Schermata di login: in basso al pannello bianco il
  pulsante \textbf{Accedi con Google} con il logo G a sinistra. Toccandolo
  si apre il selettore degli account Google presenti sul dispositivo.}{login_google.jpeg}{%
  Finestra di sistema \textbf{Choose an account}: mostra tutti gli account
  Google salvati sul dispositivo con foto profilo e indirizzo email (oscurati
  per la privacy). In fondo la voce \emph{Add another account} per aggiungere
  un account non ancora presente, e la nota che Google condividera' nome,
  email e foto profilo con NutriApp.}

\labelgreen{Descrizione}
Toccare \textbf{Accedi con Google} nella schermata di login per avviare
l'autenticazione rapida. Il sistema operativo apre automaticamente la finestra
di selezione account di Google, che mostra tutti gli account Gmail gia'
configurati sul dispositivo. Ogni voce nella lista riporta la foto profilo
e l'indirizzo email dell'account.

\labelgreen{Come procedere}
\begin{itemize}[leftmargin=1.4em, itemsep=2pt]
  \item \textbf{Toccare il proprio account} nell'elenco per selezionarlo e
        autorizzare NutriApp ad accedere a nome, email e foto profilo.
        L'accesso viene completato automaticamente senza ulteriori passaggi
  \item \textbf{Toccare \emph{Add another account}} se l'account desiderato
        non e' nella lista: si apre il flusso di login Google standard per
        aggiungere un nuovo account al dispositivo
\end{itemize}

\labelgreen{Passaggio successivo}
Dopo la selezione dell'account, l'applicazione autentica l'utente e apre
direttamente la schermata principale. Se e' il primo accesso con quell'account
Google, viene avviato il flusso di configurazione iniziale del profilo.

\suggerimento{L'accesso con Google e' il metodo piu' rapido: non richiede
di ricordare una password separata e si aggiorna automaticamente se si cambia
la password Google. E' consigliato per chi usa gia' Google su altri servizi.}

% ─── 1.3 ─────────────────────────────────────────────────────────────────────
\subsection{Creare un nuovo account}

\passaggio{Aprire il form di registrazione --- Step 1}

\schermata[0.40\textwidth]{registra_step1.jpeg}{Schermata \textbf{Registrazione}
  --- Step 1 di 3: l'indicatore di avanzamento in cima mostra il primo trattino
  verde (attivo) e due grigi (da completare). Al centro l'avatar segnaposto con
  l'icona fotocamera verde per aggiungere una foto opzionale. I campi del form
  sono: \textbf{Nome}, \textbf{Cognome}, \textbf{Email}, \textbf{Password} e
  \textbf{Conferma Password}. In fondo il pulsante verde \textbf{REGISTRATI}.}

\labelgreen{Descrizione}
Toccare il link verde \emph{Non hai un account? Registrati} nella schermata
di login per aprire il wizard di registrazione a tre step, visualizzati
dall'indicatore di avanzamento in cima (trattino verde = step corrente,
trattini grigi = step successivi).

\textbf{Step 1 --- Dati personali e credenziali.}
Compilare i seguenti campi:

\begin{itemize}[leftmargin=1.4em, itemsep=2pt]
  \item \textbf{Nome} e \textbf{Cognome} --- il nome completo che verra'
        visualizzato nell'app e nel saluto della schermata principale
  \item \textbf{Email} --- l'indirizzo di posta che funzionera' da
        identificativo univoco dell'account; deve essere valido perche'
        verra' usato per la verifica tramite codice OTP nel passo successivo
  \item \textbf{Password} --- scegliere una password sicura di almeno 8
        caratteri; si consiglia di includere lettere maiuscole, minuscole,
        numeri e caratteri speciali
  \item \textbf{Conferma Password} --- reinserire la stessa password per
        evitare errori di digitazione; i due campi devono corrispondere
        esattamente per poter procedere
\end{itemize}

\textbf{Foto profilo (opzionale).}
Toccare l'icona fotocamera verde sovrapposta all'avatar grigio per scattare
una foto o sceglierla dalla galleria. La foto e' facoltativa e puo' essere
aggiunta o modificata in qualsiasi momento successivo dalla sezione
\emph{Il mio profilo} nelle impostazioni.

\labelgreen{Passaggio successivo}
Controllare che tutti i campi obbligatori siano compilati correttamente,
quindi premere il pulsante verde \textbf{REGISTRATI} per passare allo step
successivo.

\passaggio{Verificare l'indirizzo email --- Step 2}

\schermata[0.40\textwidth]{registra_step2.jpeg}{Schermata \textbf{Registrazione}
  --- Step 2 di 3: l'indicatore mostra i primi due trattini verdi e l'ultimo
  grigio. Al centro l'icona di una busta con spunta verde e il titolo
  \textbf{Verifica la tua email}. Il testo indica che e' stato inviato un
  codice OTP all'indirizzo email inserito (visualizzato parzialmente oscurato).
  Il campo \textbf{Codice OTP} attende l'inserimento. Il link 
  \emph{Invia di nuovo} permette di richiederne uno nuovo.
  In fondo il pulsante verde \textbf{CONTROLLA EMAIL}.}

\labelgreen{Descrizione}
Dopo aver premuto \textbf{Continua} nel primo step, l'applicazione invia
automaticamente un codice OTP (One-Time Password) all'indirizzo email
inserito. La schermata mostra l'icona di una busta con spunta verde e
indica a quale indirizzo e' stato spedito il codice.

\labelgreen{Come completare la verifica}
\begin{itemize}[leftmargin=1.4em, itemsep=2pt]
  \item Aprire la casella di posta dell'indirizzo email usato durante la
        registrazione e cercare l'email di NutriApp relativa alla verifica
  \item Annotare o copiare il codice OTP a 6 cifre presente nell'email
  \item Tornare nell'app e inserire il codice nel campo \textbf{Codice OTP}
  \item Premere \textbf{CONTROLLA EMAIL} per confermare l'indirizzo email
        e procedere all'ultimo step
\end{itemize}

\labelgreen{Codice non ricevuto}
Se il codice non arriva entro qualche minuto, verificare la cartella
\emph{Spam} o \emph{Posta indesiderata} del proprio provider. In alternativa,
toccare il link \emph{Invia di nuovo} per richiedere un nuovo codice OTP.
Il link diventa attivo dopo un breve intervallo di attesa per evitare
invii multipli ravvicinati.

\suggerimento{Il codice OTP ha una validita' limitata nel tempo (di solito
10--15 minuti). Se si attende troppo a lungo prima di inserirlo, sara'
necessario richiederne uno nuovo tramite il link \emph{Invia di nuovo}.}

\passaggio{Completare la registrazione --- Step 3}

\schermata[0.40\textwidth]{registra_step3.jpeg}{Schermata \textbf{Registrazione}
  --- Step 3 di 3: tutti e tre i trattini dell'indicatore sono verdi a segnalare
  il completamento del percorso. Al centro un cerchio verde con spunta bianca e
  il titolo \textbf{Ottimo lavoro!} con il sottotitolo \emph{Sei pronto a
  iniziare il tuo percorso con NutriApp.} In fondo il pulsante verde
  \textbf{VAI ALLA HOMEPAGE}.}

\labelgreen{Descrizione}
Dopo la verifica riuscita del codice OTP, l'app mostra la schermata di
conferma. Tutti e tre i trattini dell'indicatore diventano verdi. Il
messaggio \textbf{Ottimo lavoro!} conferma che l'account e' stato creato e
verificato correttamente. Premere \textbf{VAI ALLA HOMEPAGE} per accedere
all'applicazione.

\begin{infobox}
  \textbf{\color{nutrigreen}Riepilogo delle modalita' di accesso}

  \smallskip
  \begin{tabular}{@{}lll@{}}
    \toprule
    \textbf{Metodo}       & \textbf{Requisito}             & \textbf{Ideale per} \\
    \midrule
    Email + Password      & Account NutriApp               & Chi preferisce credenziali dedicate \\
    Accedi con Google     & Account Gmail sul dispositivo  & Accesso rapido senza password \\
    Accedi con Apple      & Apple ID (solo iOS)            & Utenti iPhone/iPad \\
    Senza accedere        & Nessuno                        & Esplorazione rapida dell'app \\
    \bottomrule
  \end{tabular}
\end{infobox}

\newpage

% ════════════════════════════════════════════════════════════════════════════
%  2. BARRA DI NAVIGAZIONE
% ════════════════════════════════════════════════════════════════════════════
\section{Barra di navigazione}

La barra di navigazione e' l'elemento persistente dell'interfaccia: appare nella
parte inferiore di ogni schermata e consente di spostarsi istantaneamente tra le
cinque aree principali dell'applicazione con un singolo tocco.

% ─── 2.1 ─────────────────────────────────────────────────────────────────────
\subsection{Visualizzare la barra di navigazione}

\passaggio{Aprire l'applicazione}

\schermata[0.38\textwidth]{homepage.jpg}{Schermata principale \textbf{NutriApp}:
  riepilogo calorico giornaliero con grafico circolare, barre dei macronutrienti
  (Carboidrati, Proteine, Grassi) e le quattro fasce pasto. La barra di
  navigazione in basso con le icone NutriApp, Calendario, Aggiungi, Ricette
  e Grafici.}

\labelgreen{Descrizione}
All'apertura dell'applicazione viene mostrata la schermata principale con il
saluto personalizzato in cima, la data corrente navigabile con le frecce
\texttt{<} e \texttt{>}, e un grafico circolare verde che indica le calorie
rimanenti nella giornata. Sulla destra del grafico sono visibili le barre
orizzontali colorate per Carboidrati (arancione), Proteine (verde) e Grassi
(rosa/rosso), ciascuna con il valore consumato rispetto all'obiettivo.
Le quattro fasce orarie --- Colazione, Pranzo, Cena e Snack --- sono
accessibili con il pulsante \textbf{+} a fianco di ciascuna.

\labelgreen{Elementi della barra}
\begin{itemize}[leftmargin=1.4em, itemsep=2pt]
  \item \textbf{NutriApp} (icona casetta) --- schermata principale con riepilogo
        calorie, macronutrienti e sezioni pasto giornaliere
  \item \textbf{Calendario} (icona calendario) --- vista mensile con codice
        colore verde/rosso per i giorni con o senza pasti registrati
  \item \textbf{Aggiungi} (pulsante \textbf{+} verde centrale) --- accesso
        rapido alla schermata di inserimento alimenti
  \item \textbf{Ricette} (icona libro) --- elenco delle ricette salvate
  \item \textbf{Grafici} (icona grafico a linee) --- statistiche nutrizionali
        con viste giornaliera, settimanale, mensile e annuale
\end{itemize}

\suggerimento{L'icona attiva viene evidenziata in verde nella barra inferiore.
Il pulsante \textbf{+} verde al centro e' sempre disponibile per aggiungere
velocemente un alimento: e' il percorso piu' rapido per la registrazione quotidiana.}

\newpage

% ════════════════════════════════════════════════════════════════════════════
%  3. INSERIMENTO DEGLI ALIMENTI
% ════════════════════════════════════════════════════════════════════════════
\section{Inserimento degli alimenti}

L'applicazione offre piu' percorsi per registrare i pasti: ricerca nel database
globale, selezione tra alimenti salvati, aggiunta di una ricetta gia' creata e
scansione del codice a barre del prodotto.

% ─── 3.1 ─────────────────────────────────────────────────────────────────────
\subsection{Aprire la schermata di inserimento}

\passaggio{Accedere alla schermata principale}

\schermata[0.38\textwidth]{homepage.jpg}{HomePage: grafico circolare con le calorie
  rimanenti, barre dei macronutrienti e le quattro fasce orarie ciascuna con
  il pulsante \textbf{+}.}

\labelgreen{Passaggio successivo}
Premere il \textbf{+} accanto alla fascia oraria desiderata (Colazione, Pranzo,
Cena o Snack) oppure premere il grande pulsante verde \textbf{+} al centro della
barra di navigazione per accedere alla schermata \emph{Aggiungi}.

% ─── 3.2 ─────────────────────────────────────────────────────────────────────
\subsection{Scegliere la modalita' di inserimento}

\passaggio{Selezionare la sorgente dell'alimento}

\dueschermate{pasto_dropdown.jpg}{Schermata \textbf{Manuale}: menu a tendina
  \emph{Seleziona Pasto} e campi Nome Alimento, Peso (g) e Macronutrienti
  (Calorie, Carboidrati, Proteine, Grassi, Fibre, Zuccheri).}{pasto_recenti.jpg}{%
  Schermata \textbf{Aggiungi} con tab \emph{Aggiungi Alimento} attivo: barra
  di ricerca con icona barcode. In basso il link \emph{Non trovi alimento?
  Registralo te!}}

\labelgreen{Descrizione}
La schermata \emph{Aggiungi} presenta tre tab nella parte superiore:

\begin{itemize}[leftmargin=1.4em, itemsep=2pt]
  \item \textbf{Aggiungi Alimento} --- ricerca nel database globale o scansione
        del codice a barre
  \item \textbf{Ricette Salvate} --- le ricette create dall'utente con il totale
        calorico di ciascuna
  \item \textbf{Alimenti Salvati} --- gli alimenti registrati manualmente in
        precedenza
\end{itemize}

\passaggio{Navigare tra i tab Ricette Salvate e Alimenti Salvati}

\dueschermate{pasto_ricette.jpg}{Tab \textbf{Ricette Salvate}: ogni ricetta
  mostra nome e calorie totali. In fondo il link
  \emph{Inserisci nuova ricetta}.}{pasto_preferiti.jpg}{Tab \textbf{Alimenti
  Salvati}: barra di ricerca; se non ci sono alimenti salvati compare il
  messaggio \emph{Nessun alimento salvato}.}

\suggerimento{Gli alimenti inseriti manualmente vengono salvati automaticamente
nella sezione \emph{Alimenti Salvati} e saranno disponibili per gli inserimenti
futuri senza doverli reinserire da capo.}

% ─── 3.3 ─────────────────────────────────────────────────────────────────────
\subsection{Inserimento manuale di un alimento}

\passaggio{Compilare il form di inserimento}

\dueschermate{registra_alimento.jpg}{Form \textbf{Ingrediente}: campi
  \emph{Nome Alimento} e \emph{Peso (g)}, seguiti dalla sezione
  \textbf{Macronutrienti} con Carboidrati, Proteine, Grassi, Fibre,
  Zuccheri e Acqua.}{nutrienti_micro1.jpg}{Sezione \textbf{Ingrediente} con
  le sezioni espandibili \emph{Dettagli Grassi}, \emph{Vitamine} e
  \emph{Minerali} per un tracciamento nutrizionale completo.}

\labelgreen{Descrizione}
Toccare \emph{Registralo te!} per aprire il form di inserimento manuale.
Compilare almeno i campi obbligatori della sezione \textbf{Macronutrienti}
leggendo i valori dalla confezione del prodotto (etichetta \emph{Valori
nutrizionali per 100\,g}). Per una tracciatura completa e' possibile espandere
le sezioni \textbf{Dettagli Grassi}, \textbf{Vitamine} e \textbf{Minerali}.


% ─── 3.4 ─────────────────────────────────────────────────────────────────────
\subsection{Inserimento tramite scansione del codice a barre}

\passaggio{Avviare la fotocamera di scansione}

\schermata[0.38\textwidth]{scansione_barcode.jpg}{Schermata di scansione
  \textbf{barcode}: la fotocamera inquadra il retro di una confezione.
  Il riquadro verde delimita l'area di riconoscimento. In alto i controlli
  \textbf{FLASH ON} e \textbf{CANCEL}.}

\labelgreen{Descrizione}
Toccare l'icona barcode a destra della barra di ricerca nel tab
\emph{Aggiungi Alimento}. Inquadrare il codice a barre all'interno del
riquadro verde: il riconoscimento avviene in automatico non appena il
codice entra a fuoco, senza premere alcun pulsante.

\suggerimento{In ambienti poco illuminati attivare il flash toccando
\textbf{FLASH ON} in alto a destra. Tenere il telefono a circa 15--20\,cm
dalla confezione per un risultato ottimale.}

% ─── 3.5 ─────────────────────────────────────────────────────────────────────
\subsection{Creare una nuova ricetta}

\passaggio{Aprire il form e aggiungere gli ingredienti}

\treschermate{nuova_ricetta.jpg}{Form \textbf{Crea Ricetta}: campi
  \emph{Nome ricetta}, \emph{Porzioni} e \emph{Procedimento / Note}.
  La sezione Ingredienti mostra \emph{0 aggiunti}. In fondo i pulsanti
  \textbf{+ Ingredienti} e \textbf{Salva}.}{nuova_ricetta_recenti.jpg}{%
  Schermata \textbf{Aggiungi Ingrediente} con tab \emph{Aggiungi Alimento}:
  barra di ricerca con barcode. In basso il link
  \emph{Non trovi alimento? Registralo te!}}{nuova_ricetta_preferiti.jpg}{%
  Tab \emph{Alimenti Salvati}: elenco degli alimenti gia' registrati o
  messaggio \emph{Nessun alimento salvato} se vuoto.}

\labelgreen{Descrizione}
Dalla sezione \emph{Ricette} toccare il pulsante \textbf{+} verde per aprire
il form \emph{Crea Ricetta}. Compilare nome, numero di porzioni e note
opzionali, poi premere \textbf{+ Ingredienti} per aggiungere ciascun
componente. Ogni ingrediente aggiunto appare nella lista con le icone di
modifica e cancellazione. Al termine premere \textbf{Salva}.

% ─── 3.6 ─────────────────────────────────────────────────────────────────────
\subsection{Consultare il dettaglio di un pasto}

\passaggio{Aprire la schermata dettaglio fascia oraria}

\schermata[0.38\textwidth]{dettaglia_pasto.jpg}{Schermata \textbf{Colazione}: banner arancione
  con icona e contatore alimenti, navigazione per data, pulsante \textbf{+}
  per aggiungere alimenti. Il \emph{Riepilogo Nutrizionale} mostra barre di
  avanzamento per Calorie, Carboidrati, Proteine e Grassi. In fondo la lista
  \emph{Alimenti Consumati} con icone modifica e cancellazione.}

\labelgreen{Descrizione}
Toccare una fascia oraria nella schermata principale per aprire il riepilogo
dettagliato con il banner colorato, il riepilogo nutrizionale con barre di
avanzamento e la lista degli alimenti consumati con le icone di modifica
(matita blu) e cancellazione (cestino rosso).

\newpage

% ════════════════════════════════════════════════════════════════════════════
%  4. IMPOSTAZIONE DEGLI OBIETTIVI
% ════════════════════════════════════════════════════════════════════════════
\section{Impostazione degli obiettivi}

Gli obiettivi definiscono i valori di riferimento per calorie, macronutrienti
e peso che l'app usa per calcolare i progressi. Impostare obiettivi realistici
e' fondamentale per ottenere dati di avanzamento significativi.

% ─── 4.1 ─────────────────────────────────────────────────────────────────────
\subsection{Accedere e configurare gli obiettivi}

\passaggio{Aprire le Impostazioni e selezionare I miei obiettivi}

\dueschermate{impostazioni.jpg}{Schermata \textbf{Impostazioni}: le sezioni Account
  (Il mio profilo, I miei obiettivi, Notifiche, Lingua e
  paese), Supporto (Aiuto, Informativa sulla privacy, Su di noi) e Azioni (Segnala un
  problema, Log out).}{obiettivi.jpg}{%
  Schermata \textbf{I miei obiettivi}: sezione \emph{Peso} con \emph{Peso
  attuale} (60.0\,kg) e \emph{Obiettivo peso} (70.0\,kg). Sezione
  \emph{Macronutrienti} con \emph{Goal calorie} (2000.0\,kcal),
  \emph{Carboidrati} (130.0\,g), \emph{Proteine} (45.0\,g) e
  \emph{Grassi} (36.0\,g). In fondo le sezioni espandibili
  \emph{Dettagli Macro}, \emph{Vitamine} e \emph{Minerali}.}

\labelgreen{Descrizione}
Toccare l'icona ingranaggio in alto a destra per aprire le impostazioni,
poi selezionare \textbf{I miei obiettivi}. La schermata raccoglie i parametri
in due sezioni:

\begin{itemize}[leftmargin=1.4em, itemsep=2pt]
  \item \textbf{Peso} --- il \emph{Peso attuale} e' il punto di partenza per
        il calcolo del fabbisogno; l'\emph{Obiettivo peso} e' il traguardo
        da raggiungere (puo' essere maggiore o minore del peso attuale)
  \item \textbf{Macronutrienti} --- il \emph{Goal calorie} definisce il limite
        calorico giornaliero; i campi Carboidrati, Proteine e Grassi impostano
        i limiti in grammi
\end{itemize}

\labelgreen{Passaggio successivo}
Inserire i valori desiderati e premere \textbf{Salva}. Le modifiche vengono
applicate immediatamente a tutta l'applicazione. Il banner verde di conferma
appare in fondo allo schermo.

\begin{infobox}
  \textbf{\color{nutrigreen}Come scegliere i valori giusti:}
  un punto di partenza comune e' carboidrati 45--65\%, proteine 10--35\%,
  grassi 20--35\% delle calorie totali giornaliere. Consultare un
  professionista della nutrizione per valori personalizzati.
\end{infobox}

\newpage

% ════════════════════════════════════════════════════════════════════════════
%  5. VISUALIZZAZIONE DEL PROGRESSO
% ════════════════════════════════════════════════════════════════════════════
\section{Visualizzazione del progresso}

L'applicazione offre quattro viste temporali per analizzare l'andamento
calorico e nutrizionale: giornaliera, settimanale, mensile e annuale,
accessibili dalla sezione \textbf{Grafici} nella barra di navigazione.

% ─── 5.1 ─────────────────────────────────────────────────────────────────────
\subsection{Accedere ai dati di progresso}

\passaggio{Aprire la sezione Grafici}

\schermata[0.38\textwidth]{storico_giornaliero.jpg}{Schermata \textbf{Grafici}
  con tab \emph{Giornaliero} attivo: riquadro verde \emph{Statistiche
  nutrizionali} con voci registrate, giorni tracciati e data dell'ultimo
  inserimento. Sotto il selettore filtri, i quattro tab temporali e il box
  rosso con il totale calorico.}

\labelgreen{Descrizione}
Toccare l'icona \textbf{Grafici} nella barra di navigazione. In cima compare
il riquadro verde con il riepilogo globale, poi il pannello \textbf{Filtri}
con il nutriente selezionato (default: Calorie) e l'intervallo di date.
I quattro tab permettono di cambiare la granularita' della visualizzazione.

% ─── 5.2 ─────────────────────────────────────────────────────────────────────
\subsection{Analizzare i grafici}

\passaggio{Grafico a barre --- Calorie nel tempo}

\schermata[0.40\textwidth]{grafici_barre.jpeg}{Schermata \textbf{Grafici} con
  tab \emph{Giornaliero} attivo. In cima il pannello \textbf{Esporta report PDF}
  con i pulsanti \textbf{Stampa PDF} (verde) e \textbf{Condividi PDF}. Il
  grafico \textbf{Calorie nel tempo} mostra tre barre verticali rosse per i
  giorni 28, 29 e 30 aprile: 779\,kcal, 143\,kcal e 377\,kcal. La scala
  sull'asse verticale va da 0.0 a 779\,kcal. In basso inizia la sezione
  \emph{Distribuzione macronutrienti}.}

\labelgreen{Descrizione}
Il grafico a barre \textbf{Calorie nel tempo} mostra l'apporto calorico
giornaliero nel periodo selezionato. Ogni barra verticale rossa rappresenta
un giorno, con il valore in kcal indicato sopra. L'altezza e' proporzionale
al totale calorico: giorni con pochi alimenti registrati producono barre
piu' basse (come il 29 aprile con 143\,kcal), giorni con piu' pasti producono
barre piu' alte (come il 28 aprile con 779\,kcal). L'asse verticale si adatta
automaticamente al valore massimo del periodo.

\labelgreen{Esporta report PDF}
Sopra il grafico il pannello \textbf{Esporta report PDF} offre due pulsanti:
\begin{itemize}[leftmargin=1.4em, itemsep=2pt]
  \item \textbf{Stampa PDF} (verde) --- genera il documento e apre la finestra
        di stampa del sistema operativo
  \item \textbf{Condividi PDF} (bianco con bordo) --- genera il documento e
        apre il pannello di condivisione per inviarlo via email, WhatsApp o
        qualsiasi altra app installata
\end{itemize}

\passaggio{Grafico a torta --- Distribuzione macronutrienti}

\schermata[0.40\textwidth]{grafici_torta.jpeg}{Schermata \textbf{Grafici} ---
  sezione \textbf{Distribuzione macronutrienti}: grafico ad anello con la
  ripartizione percentuale dei macronutrienti. La legenda indica Carboidrati
  44.7\%, Proteine 12.0\% e Grassi 43.3\%. Le barre orizzontali sotto
  mostrano i valori assoluti: Carboidrati 155.6\,g (barra arancione),
  Proteine 41.9\,g (barra blu) e Grassi 67.1\,g (barra viola).}

\labelgreen{Descrizione}
Scorrendo verso il basso nella schermata \emph{Grafici} compare la sezione
\textbf{Distribuzione macronutrienti}, con due rappresentazioni complementari:

\begin{itemize}[leftmargin=1.4em, itemsep=2pt]
  \item \textbf{Grafico ad anello} --- mostra visivamente la proporzione tra
        i tre macronutrienti: il segmento dorato/arancione rappresenta i
        Carboidrati, quello blu scuro le Proteine e quello viola i Grassi.
        Le percentuali nella legenda a destra consentono di leggere i valori
        esatti a colpo d'occhio
  \item \textbf{Barre orizzontali} --- mostrano sia la percentuale sia il
        valore assoluto in grammi di ciascun macronutriente, facilitando il
        confronto con gli obiettivi impostati nella sezione \emph{I miei obiettivi}
\end{itemize}

\labelgreen{Come interpretare i dati}
Nell'esempio mostrato, i Grassi costituiscono il 43.3\% dell'apporto
energetico totale (67.1\,g), un valore superiore alla soglia consigliata
del 20--35\%. I Carboidrati si attestano al 44.7\% (155.6\,g) e le Proteine
al 12.0\% (41.9\,g). Queste informazioni permettono di identificare squilibri
nella distribuzione dei macronutrienti e di aggiustare le scelte alimentari
nei giorni successivi.

\labelgreen{Passaggio successivo}
Usare il pannello \textbf{Filtri} in cima alla schermata per cambiare il
nutriente visualizzato (es. da Calorie a Carboidrati) e l'intervallo di date.
Toccare i tab \textbf{Settimanale}, \textbf{Mensile} o \textbf{Annuale} per
visualizzare l'andamento su periodi piu' lunghi.

\suggerimento{Il confronto tra il grafico ad anello e gli obiettivi in
\emph{I miei obiettivi} e' il modo piu' efficace per valutare se la propria
dieta e' in linea con i target. Se una fetta e' molto piu' grande delle altre,
conviene ribilanciare i pasti dei giorni successivi.}

% ─── 5.3 ─────────────────────────────────────────────────────────────────────
\subsection{Usare il Calendario}

\passaggio{Aprire il Calendario}

\schermata[0.38\textwidth]{calendario.jpg}{Schermata \textbf{Calendario} di
  Aprile 2026: giorni in rosso senza pasti registrati, il 28 in verde con
  bordo blu (giorno corrente con pasto registrato), giorni futuri in grigio.
  In basso la legenda: verde = Pasti registrati, rosso = Senza pasti.}

\labelgreen{Descrizione}
Toccare l'icona \textbf{Calendario} nella barra di navigazione. Il codice
colore identifica immediatamente i giorni registrati: verde (pasto presente),
rosso (nessun pasto), grigio chiaro (giorni futuri), bordo blu (giorno corrente).
Toccare un giorno per visualizzare il dettaglio dei pasti di quella data.

\newpage

% ════════════════════════════════════════════════════════════════════════════
%  6. IMPOSTAZIONI DELL'ACCOUNT
% ════════════════════════════════════════════════════════════════════════════
\section{Impostazioni dell'account}

La sezione \emph{Settings} raccoglie tutte le opzioni di personalizzazione:
profilo personale, notifiche, unita' di misura, lingua e gestione dell'account.

% ─── 6.1 ─────────────────────────────────────────────────────────────────────
\subsection{Aprire le impostazioni}

\passaggio{Accedere alla schermata Settings}

\schermata[0.38\textwidth]{impostazioni.jpg}{Schermata \textbf{Impostazioni}:
  le tre sezioni Account, Supporto e Azioni con tutte le voci
  disponibili, ciascuna con la propria icona.}

\labelgreen{Descrizione}
Toccare l'icona ingranaggio in alto a destra per aprire le impostazioni.

\medskip
\begin{infobox}
  \begin{tabular}{@{}ll@{}}
    \toprule
    \textbf{Voce}              & \textbf{Funzione} \\
    \midrule
    Il mio profilo             & Modifica nome, cognome, altezza, data nascita, sesso e foto \\
    I miei obiettivi           & Imposta calorie target, peso, macro e micronutrienti \\
    Notifiche                  & Promemoria pasti e idratazione con orari personalizzati \\
    Lingua e paese             & Lingua (IT, EN, ZH, AR) e paese di riferimento \\
    \midrule
    Aiuto                      & Centro assistenza con domande frequenti \\
    Informativa sulla privacy  & Documento sulle policy di gestione dati \\
    Su di noi                  & Informazioni sull'applicazione e il team \\
    \midrule
    Segnala un problema        & Segnalazione diretta al team di supporto \\
    Log out                    & Disconnessione dall'account corrente \\
    \bottomrule
  \end{tabular}
\end{infobox}

% ─── 6.2 ─────────────────────────────────────────────────────────────────────
\subsection{Modificare i dati dell'account}

\passaggio{Aggiornare il profilo personale}

\schermata[0.38\textwidth]{profilo.jpg}{Schermata \textbf{Il mio profilo}:
  foto profilo con icona fotocamera verde, campi \emph{Nome} (Christian),
  \emph{Cognome} (Donzelli), \emph{Altezza} (175.0), \emph{Data di nascita}
  (30/11/-1) e menu a tendina \emph{Sesso} (Maschio). In fondo il
  pulsante verde \textbf{Salva}.}

\labelgreen{Descrizione}
Toccare \textbf{Il mio profilo} per modificare i dati personali:
\begin{itemize}[leftmargin=1.4em, itemsep=2pt]
  \item \textbf{Foto profilo} --- toccare l'icona fotocamera verde per scattare
        una foto o sceglierne una dalla galleria
  \item \textbf{Nome e Cognome} --- le informazioni anagrafiche dell'utente
  \item \textbf{Altezza} --- in centimetri, usata per il calcolo del fabbisogno
  \item \textbf{Data di nascita} --- usata per il calcolo del metabolismo basale
  \item \textbf{Sesso} --- influisce sulla formula del fabbisogno calorico
\end{itemize}

Premere \textbf{Salva} per confermare. L'app aggiorna automaticamente i calcoli
nutrizionali in base ai nuovi dati.

% ─── 6.3 ─────────────────────────────────────────────────────────────────────
\subsection{Configurare le notifiche}

\passaggio{Attivare i promemoria}

\dueschermate{notifiche.jpg}{Schermata \textbf{Notifiche} con toggle attivi:
  \emph{Promemoria pasti} con orari Colazione 7:30\,AM, Pranzo 12:00\,PM,
  Cena 7:00\,PM, Snacks 3:00\,PM; \emph{Promemoria acqua} con frequenza
  \emph{Ogni 30 minuti}.}{notifiche_orari.jpg}{Dropdown \emph{Frequenza
  promemoria} aperto: \emph{Ogni 30 minuti} (selezionata), \emph{Ogni 1 ora},
  \emph{Ogni 2 ore} e \emph{DECIDO IO}. In fondo la \emph{Cronologia avvisi}
  con le notifiche recenti.}

\labelgreen{Descrizione}
Nella schermata \emph{Notifiche} sono disponibili due toggle:
\begin{itemize}[leftmargin=1.4em, itemsep=2pt]
  \item \textbf{Promemoria pasti} --- attivandolo compaiono le quattro fasce
        orarie; toccare un orario verde per modificarlo con il selettore di sistema
  \item \textbf{Promemoria acqua} --- attivandolo compare il menu frequenza
        (Ogni 30 min, Ogni 1 ora, Ogni 2 ore, DECIDO IO)
\end{itemize}
In fondo la \textbf{Cronologia avvisi} con le notifiche recenti; toccare
\textbf{Pulisci} per cancellarla.

% ─── 6.4 ─────────────────────────────────────────────────────────────────────
\subsection{Modificare lingua e paese}

\passaggio{Configurare localizzazione}

\schermata[0.38\textwidth]{lingua_paese.jpg}{Schermata
  \textbf{Lingua e Paese}: menu \emph{Lingua} aperto con \emph{Italiano}
  selezionato, \emph{Inglese}, \emph{Cinese semplice} e \emph{Arabo}.
  Sotto la sezione \emph{Paese} con \emph{Italia}, \emph{Inghilterra},
  \emph{Cina} ed \emph{Egitto}.}

\labelgreen{Lingua e paese}
Le lingue disponibili sono Italiano, Inglese, Cinese semplice e Arabo.
La modifica si applica all'interfaccia al riavvio dell'app.

% ─── 6.5 ─────────────────────────────────────────────────────────────────────
\subsection{Segnalare un problema}

\passaggio{Inviare una segnalazione al supporto}

\schermata[0.38\textwidth]{segnala_problema.jpg}{Schermata \textbf{Segnala un
  problema}: campo di testo con placeholder \emph{Descrivi dettagliatamente
  il problema riscontrato\ldots} e pulsante verde \textbf{INVIA SEGNALAZIONE}.}

\labelgreen{Descrizione}
Toccare \textbf{Segnala un problema} nella sezione \emph{Actions}. Compilare il
campo di testo con una descrizione dettagliata e premere \textbf{INVIA SEGNALAZIONE}.

\suggerimento{Per una risoluzione piu' rapida, includere nella segnalazione il
modello del dispositivo (es.\ iPhone 14, Samsung Galaxy S23), la versione del
sistema operativo e la descrizione precisa dei passaggi che hanno portato al problema.}

\bigskip
\noindent\rule{\textwidth}{1pt}{\color{nutrigreen}}
\begin{center}
  \small\itshape\color{nutrigray}--- Fine Sezione 4 ---
\end{center}

\end{document}
